\documentclass[11pt,a4paper]{article}
\usepackage[utf8]{inputenc}
\usepackage[T1]{fontenc}
\usepackage{amsmath}
\usepackage{amssymb}
\usepackage{amsthm}
\usepackage{enumitem}
\usepackage{hyperref}
\usepackage{geometry}
\geometry{margin=1in}

\title{Mathematical Interpretation of test\_encode.c Using Coding Theory}
\author{}
\date{}

\begin{document}

\maketitle

\section{Overview}

The unit tests in \texttt{test\_encode.c} verify a bit-level encoding function \texttt{dap\_encode\_char\_by\_char} that implements a \textbf{base-2\textsuperscript{b} encoding scheme}, where \texttt{b = a\_base\_size}. From a coding theory perspective, this function implements a \textbf{linear encoding map} from a binary vector space to a representation space.

\section{Mathematical Framework}

\subsection{Vector Space Structure}

Let $\mathbb{F}_2 = \{0, 1\}$ denote the binary field (Boolean field). The input data is interpreted as a vector in the vector space:

$$\mathbf{x} \in \mathbb{F}_2^n$$

where $n = 8 \cdot \text{input\_size}$ is the total number of bits in the input.

\subsection{Encoding Map}

The function implements an encoding map:

$$E: \mathbb{F}_2^n \rightarrow \Sigma^m$$

where:
\begin{itemize}
\item $\Sigma$ is the output alphabet (defined by the lookup table)
\item $m = \lfloor n/b \rfloor$ is the output length
\item $b = \text{a\_base\_size}$ is the number of bits extracted per output symbol
\end{itemize}

\subsection{Code Definition}

The encoding process defines a \textbf{code} $C \subseteq \mathbb{F}_2^n$ where each codeword $\mathbf{c} \in C$ represents an input vector. However, since this is a \textbf{bijective encoding} (not error-correcting), the code $C = \mathbb{F}_2^n$ (the entire space).

The code can be viewed as a \textbf{linear code} with:
\begin{itemize}
\item \textbf{Dimension}: $k = n$ (all input vectors are valid)
\item \textbf{Length}: $n$ (input length in bits)
\item \textbf{Rate}: $R = 1$ (no redundancy added)
\end{itemize}

\section{Test-by-Test Mathematical Interpretation}

\subsection{1. \texttt{s\_test\_encode\_null\_inputs()} --- Trivial Cases}

\textbf{Mathematical Interpretation:}
\begin{itemize}
\item Tests the \textbf{domain} of the encoding function $E$
\item Verifies that $E(\emptyset) = \emptyset$ (empty input maps to empty output)
\item Ensures the function is \textbf{well-defined} on its domain
\end{itemize}

\textbf{Coding Theory Perspective:}
\begin{itemize}
\item Verifies that the encoding map respects the \textbf{zero vector}: $E(\mathbf{0}) = \emptyset$
\item Tests \textbf{function domain constraints}: $E$ is only defined for non-null inputs
\end{itemize}

\subsection{2. \texttt{s\_test\_encode\_empty\_input()} --- Zero Vector}

\textbf{Mathematical Interpretation:}
\begin{itemize}
\item Tests encoding of the \textbf{zero vector} $\mathbf{0} \in \mathbb{F}_2^0$
\item Verifies: $E(\mathbf{0}) = \emptyset$ (empty string)
\end{itemize}

\textbf{Coding Theory Perspective:}
\begin{itemize}
\item The zero vector is always a codeword in any linear code
\item Verifies the \textbf{additive identity property}: $E(\mathbf{0}) = \mathbf{0}$
\end{itemize}

\subsection{3. \texttt{s\_test\_encode\_base8()} --- Identity-like Encoding}

\textbf{Mathematical Interpretation:}
\begin{itemize}
\item Tests encoding with $b = 8$ (extracts 8 bits at a time)
\item Input: $\mathbf{x} \in \mathbb{F}_2^{40}$ (5 bytes = 40 bits for ``Hello'')
\item Output length: $m = \lfloor 40/8 \rfloor = 5$
\end{itemize}

\textbf{Coding Theory Perspective:}
\begin{itemize}
\item For $b = 8$, the encoding extracts \textbf{one byte} at a time
\item The encoding map becomes: $E: \mathbb{F}_2^{8k} \rightarrow \Sigma^k$ where $k$ is the number of bytes
\item This is essentially a \textbf{byte-to-symbol mapping} with no bit-level transformation
\end{itemize}

\textbf{Vector Space View:}
\begin{itemize}
\item Each 8-bit block $\mathbf{x}_i \in \mathbb{F}_2^8$ is mapped to a symbol via the lookup table
\item The encoding is: $E(\mathbf{x}) = (T(x_1), T(x_2), \ldots, T(x_k))$ where $T$ is the table mapping
\end{itemize}

\subsection{4. \texttt{s\_test\_encode\_base5()} --- Base-32-like Encoding}

\textbf{Mathematical Interpretation:}
\begin{itemize}
\item Tests encoding with $b = 5$ (extracts 5 bits at a time)
\item Input: $\mathbf{x} \in \mathbb{F}_2^{16}$ (2 bytes = 16 bits for ``AB'')
\item Output length: $m = \lfloor 16/5 \rfloor = 3$
\end{itemize}

\textbf{Coding Theory Perspective:}
\begin{itemize}
\item The encoding extracts \textbf{5-bit blocks} from the input bitstream
\item Each 5-bit block $\mathbf{b}_i \in \mathbb{F}_2^5$ represents a value in $\{0, 1, \ldots, 31\}$
\item The encoding map: $E: \mathbb{F}_2^{16} \rightarrow \Sigma^3$
\end{itemize}

\textbf{Vector Space Decomposition:}
\begin{itemize}
\item The input space $\mathbb{F}_2^{16}$ is partitioned into 3 blocks (with the last block potentially incomplete)
\item Each block spans across byte boundaries, requiring bit-level extraction
\end{itemize}

\subsection{5. \texttt{s\_test\_encode\_base6()} --- Base-64-like Encoding}

\textbf{Mathematical Interpretation:}
\begin{itemize}
\item Tests encoding with $b = 6$ (extracts 6 bits at a time)
\item Input: $\mathbf{x} \in \mathbb{F}_2^{24}$ (3 bytes = 24 bits for ``ABC'')
\item Output length: $m = \lfloor 24/6 \rfloor = 4$
\end{itemize}

\textbf{Coding Theory Perspective:}
\begin{itemize}
\item This is the standard \textbf{Base64 encoding scheme}
\item Each 6-bit block $\mathbf{b}_i \in \mathbb{F}_2^6$ represents a value in $\{0, 1, \ldots, 63\}$
\item The encoding map: $E: \mathbb{F}_2^{24} \rightarrow \Sigma^4$
\end{itemize}

\textbf{Mathematical Structure:}
\begin{itemize}
\item The input is divided into \textbf{4 blocks of 6 bits each}
\item Perfect alignment: $24 = 4 \times 6$ (no padding needed)
\item This is why Base64 is commonly used: 3 bytes map cleanly to 4 characters
\end{itemize}

\subsection{6. \texttt{s\_test\_encode\_different\_sizes()} --- Dimension Testing}

\textbf{Mathematical Interpretation:}
\begin{itemize}
\item Tests the encoding function across different \textbf{input dimensions}
\item Verifies the \textbf{dimension formula}: $m = \lfloor n/b \rfloor$
\end{itemize}

\textbf{Test Cases:}
\begin{enumerate}
\item $n = 8$ bits, $b = 5$: $m = \lfloor 8/5 \rfloor = 1$
\item $n = 32$ bits, $b = 5$: $m = \lfloor 32/5 \rfloor = 6$
\item $n = 64$ bits, $b = 5$: $m = \lfloor 64/5 \rfloor = 12$
\end{enumerate}

\textbf{Coding Theory Perspective:}
\begin{itemize}
\item Verifies that the encoding map $E: \mathbb{F}_2^n \rightarrow \Sigma^m$ has the correct \textbf{codomain dimension}
\item Tests the \textbf{linearity property}: the output dimension scales correctly with input dimension
\end{itemize}

\subsection{7. \texttt{s\_test\_encode\_output\_size()} --- Dimension Formula Verification}

\textbf{Mathematical Interpretation:}
\begin{itemize}
\item Systematically tests the \textbf{dimension relationship}: $m = \lfloor n/b \rfloor$
\item Tests multiple values of $b \in \{1, 2, 4, 8\}$
\end{itemize}

\textbf{Coding Theory Perspective:}
\begin{itemize}
\item For a linear code, the \textbf{rate} is $R = k/n$ where $k$ is the dimension
\item Here, $k = n$ (no redundancy), so $R = 1$
\item The output dimension $m$ represents the \textbf{representation length} in the new alphabet
\end{itemize}

\textbf{Mathematical Verification:}
\begin{itemize}
\item $b = 1$: $m = n$ (bit-by-bit encoding)
\item $b = 2$: $m = n/2$ (2 bits per symbol)
\item $b = 4$: $m = n/4$ (4 bits per symbol, like hexadecimal)
\item $b = 8$: $m = n/8$ (byte-by-byte encoding)
\end{itemize}

\subsection{8. \texttt{s\_test\_encode\_custom\_table()} --- Encoding Function Verification}

\textbf{Mathematical Interpretation:}
\begin{itemize}
\item Tests that the encoding uses the \textbf{lookup table} $T: \{0, 1, \ldots, 2^b-1\} \rightarrow \Sigma$
\item Verifies the \textbf{encoding function composition}: $E(\mathbf{x}) = (T(b_1), T(b_2), \ldots, T(b_m))$
\end{itemize}

\textbf{Coding Theory Perspective:}
\begin{itemize}
\item The lookup table defines the \textbf{symbol mapping} for the code
\item Different tables create different \textbf{representations} of the same underlying code
\item The code structure (vector space) remains the same, only the \textbf{presentation} changes
\end{itemize}

\subsection{9. \texttt{s\_test\_encode\_base58()} --- Base-58 Encoding}

\textbf{Mathematical Interpretation:}
\begin{itemize}
\item Tests Base-58 encoding with $b = 6$ (uses 6-bit extraction, but only 58 symbols)
\item Input: $\mathbf{x} \in \mathbb{F}_2^{40}$ (5 bytes = 40 bits)
\item Output length: $m = \lfloor 40/6 \rfloor = 6$
\end{itemize}

\textbf{Coding Theory Perspective:}
\begin{itemize}
\item Base-58 is a \textbf{non-standard base} encoding (58 symbols instead of 64)
\item The encoding extracts 6 bits (values 0--63) but maps only 58 of them to valid symbols
\item This creates a \textbf{subset code} where some 6-bit patterns are invalid in the output
\end{itemize}

\textbf{Mathematical Structure:}
\begin{itemize}
\item The code $C \subseteq \mathbb{F}_2^n$ remains the same
\item The encoding map $E: \mathbb{F}_2^n \rightarrow \Sigma_{58}^m$ uses a \textbf{58-symbol alphabet}
\item The lookup table $T: \{0, 1, \ldots, 63\} \rightarrow \Sigma_{58} \cup \{\text{invalid}\}$ is a \textbf{partial function}
\end{itemize}

\subsection{10. \texttt{s\_test\_encode\_base64\_standard()} and \texttt{s\_test\_encode\_base64\_url\_safe()} --- Base-64 Variants}

\textbf{Mathematical Interpretation:}
\begin{itemize}
\item Both test Base-64 encoding with $b = 6$
\item They differ only in the \textbf{lookup table} (symbol mapping)
\end{itemize}

\textbf{Coding Theory Perspective:}
\begin{itemize}
\item Both implement the \textbf{same code} $C = \mathbb{F}_2^n$
\item They differ in the \textbf{presentation}: standard Base64 vs.\ URL-safe Base64
\item The encoding maps are \textbf{equivalent} up to symbol renaming
\end{itemize}

\textbf{Mathematical Equivalence:}
\begin{itemize}
\item Standard: $E_{\text{std}}: \mathbb{F}_2^n \rightarrow \Sigma_{\text{std}}^m$
\item URL-safe: $E_{\text{url}}: \mathbb{F}_2^n \rightarrow \Sigma_{\text{url}}^m$
\item There exists a \textbf{bijection} $\phi: \Sigma_{\text{std}} \rightarrow \Sigma_{\text{url}}$ such that $E_{\text{url}} = \phi \circ E_{\text{std}}$
\end{itemize}

\section{General Mathematical Properties}

\subsection{Linearity (Partial)}

The encoding function is \textbf{not fully linear} in the algebraic sense because:
\begin{enumerate}
\item It operates on \textbf{bit-level extraction} across byte boundaries
\item The lookup table introduces \textbf{non-linear symbol mapping}
\end{enumerate}

However, the \textbf{bit extraction} process is linear:
\begin{itemize}
\item The bit extraction: $\mathbf{b}_i = \text{extract}_b(\mathbf{x}, i)$ is a \textbf{linear operation}
\item The table lookup: $T(\mathbf{b}_i)$ is a \textbf{non-linear symbol mapping}
\end{itemize}

\subsection{Vector Space Structure}

The input space maintains the structure:
\begin{itemize}
\item \textbf{Addition}: Bitwise XOR in $\mathbb{F}_2^n$
\item \textbf{Scalar multiplication}: Multiplication in $\mathbb{F}_2$ (trivial: $0 \cdot \mathbf{x} = \mathbf{0}$, $1 \cdot \mathbf{x} = \mathbf{x}$)
\end{itemize}

\subsection{Code Properties}

\begin{itemize}
\item \textbf{Code dimension}: $k = n$ (all input vectors are valid codewords)
\item \textbf{Code length}: $n$ (input length in bits)
\item \textbf{Minimum distance}: Not applicable (this is not an error-correcting code)
\item \textbf{Rate}: $R = 1$ (no redundancy)
\end{itemize}

\subsection{Encoding Efficiency}

The encoding efficiency is determined by:
\begin{itemize}
\item \textbf{Input dimension}: $n$ bits
\item \textbf{Output dimension}: $m = \lfloor n/b \rfloor$ symbols
\item \textbf{Alphabet size}: $|\Sigma| = 2^b$ (for Base-64: 64 symbols)
\end{itemize}

The \textbf{expansion factor} is: $\frac{m}{n} = \frac{1}{b}$ symbols per bit

\section{Conclusion}

The unit tests verify a \textbf{bit-level encoding scheme} that implements a mapping from binary vectors $\mathbb{F}_2^n$ to symbol strings in an alphabet $\Sigma$. While not a traditional error-correcting code, the encoding maintains the vector space structure of the input and provides a systematic way to represent binary data in different bases (Base-5, Base-6/Base-64, Base-58, etc.).

The tests verify:
\begin{enumerate}
\item \textbf{Domain and codomain} correctness
\item \textbf{Dimension relationships} between input and output
\item \textbf{Encoding function} correctness for various base sizes
\item \textbf{Symbol mapping} via lookup tables
\item \textbf{Edge cases} (empty input, null inputs)
\end{enumerate}

From a coding theory perspective, this is a \textbf{bijective encoding} (no redundancy) that preserves all information while changing the representation format.

\end{document}
